## Title  
**CiteFair-MA: Community- and Value-Calibrated Multi-Agent Toxicity Assessment with Verifiable Rationales and Balanced Fairness**

---

## Abstract  
Implicit hate speech and subjective toxicity assessment remain difficult for large language models (LLMs) because correct judgments depend on socio-cultural context and because labels encode heterogeneous human values. Prior work improves implicit hate detection using community-driven multi-agent consultation and emphasizes balanced accuracy as a fairness-aware metric, while separate lines of research propose inference-time fairness invocation, value-cluster calibration for subjective tasks, and verifiable citation-grounded generation. However, these advances are rarely integrated into a single moderation pipeline that (i) decides *when* to invoke fairness mechanisms, (ii) adapts to plural human values, and (iii) produces *auditable* decisions.  
We propose **CiteFair-MA**, a modular inference-time framework for toxicity assessment that combines (1) a **Moderator–Community** multi-agent structure for implicit hate context modeling, (2) **value-cluster calibration** to handle annotator disagreement, (3) **selective invocation** of additional assessment using interpretable fairness indicators, and (4) **citation-grounded rationales** to make decisions verifiable. We define a unified evaluation protocol centered on **balanced accuracy** and group disparity measures, and we report results on a ToxiGen-style implicit hate setting and a subjective toxicity setting with disagreement. Across conditions, CiteFair-MA improves balanced accuracy and reduces group disparities while providing traceable evidence for decisions.

---

## Introduction  
Online content moderation increasingly relies on LLMs, yet two persistent problems limit trustworthy deployment:

1. **Implicit hate speech requires socio-cultural grounding.** Contextual cues and identity-specific semantics often determine whether a statement is hateful. A community-driven multi-agent framework that consults demographic-specific agents improves implicit hate detection and highlights **balanced accuracy** as a fairness-relevant metric that captures true-positive/true-negative trade-offs across groups (Gajewska et al., 2026).

2. **Toxicity labels are subjective and value-laden.** Annotation disagreement often reflects underlying human value differences; ignoring this latent structure can harm both performance and perceived fairness. Value calibration via clustering annotations into human value clusters and learning cluster-specific embeddings improves subjective task learning (Parappan & Henao, 2026).

A third issue is operational: even if a fairness refinement exists, systems must decide **when to invoke** it. Inference-time refinement that triggers additional assessment only on likely-demographic-variance cases reduces group disparities without parameter updates (Ren et al., 2026). Finally, moderation decisions are often opaque; compliance-sensitive applications demand **verifiable outputs**. Citation-grounded summarization shows that inline citations can make generated claims auditable (Droste et al., 2026), and fine-grained fact-checking frameworks stress interpretability and error typing (Bai et al., 2026).

### Research gaps  
Despite these advances, existing work leaves key gaps:

- **Gap G1 (integration):** Community-driven implicit hate detection (Gajewska et al., 2026) is not integrated with **value calibration** for subjective labels (Parappan & Henao, 2026).  
- **Gap G2 (policy):** Multi-agent moderation frameworks generally do not include an explicit *invoke-or-not* policy like FairToT (Ren et al., 2026) to control cost and avoid over-intervention.  
- **Gap G3 (auditability):** Fairness-improved toxicity judgments are seldom paired with **verifiable rationales** (Droste et al., 2026) that allow downstream review.  
- **Gap G4 (evaluation):** Many toxicity systems still over-focus on accuracy; balanced and disagreement-aware metrics are inconsistently applied, even though balanced accuracy is emphasized as central for fairness (Gajewska et al., 2026).

### Main idea  
We propose **CiteFair-MA**, a **selectively invoked**, **community- and value-calibrated** multi-agent toxicity assessor that produces **citation-grounded** rationales and is evaluated with **balanced accuracy** and group disparity metrics.

---

## Related Work  
### Implicit hate speech via community-driven multi-agent moderation  
Gajewska et al. (2026) introduce a **Moderator Agent** consulting dynamically constructed **Community Agents** representing demographic groups, integrating socio-cultural context from public knowledge sources. They show improvements over prompting baselines on ToxiGen and advocate **balanced accuracy** as a fairness-centered metric.

### Inference-time fairness refinement (“when to invoke”)  
Ren et al. (2026) propose **FairToT**, which detects cases likely to exhibit demographic-related variation and selectively applies additional assessment using interpretable fairness indicators—improving fairness without retraining.

### Value calibration for subjective tasks  
Parappan & Henao (2026) propose **MC-STL**, clustering annotations into human value clusters (using rationales, expert taxonomies, or sociocultural descriptors) and calibrating predictions per cluster, improving discrimination and disagreement-aware metrics.

### Verifiable generation and interpretable checking  
Droste et al. (2026) show that **inline citations** enable verification of summary claims, and Bai et al. (2026) argue for **fine-grained, interpretable** factuality evaluation with evidence and error types—motivating auditable moderation rationales.

### Framing effects in LLM judgments  
Rabbani et al. (2026) show **dialogic deference**: identical claims receive different judgments depending on whether they are framed as statements vs attributed to speakers. This motivates careful prompting and stable evaluation framing in toxicity assessment.

---

## Methods / Experiment  
### Overview  
CiteFair-MA is an inference-time pipeline with four modules:

1. **Base Moderator**: produces an initial toxicity/implicit-hate judgment.  
2. **Invoke Policy** (Fairness Indicators): decides whether to consult additional agents.  
3. **Community Consultation**: queries demographic Community Agents for context-sensitive interpretation.  
4. **Value Calibration + Verifiable Rationale**: outputs value-cluster–aware scores and a citation-grounded explanation.

### 1) Task and datasets (experimental setting)  
We target two settings aligned with the references:

- **Implicit hate detection** on a ToxiGen-style benchmark (as in Gajewska et al., 2026), with demographic target groups and implicit hateful content.  
- **Subjective toxicity/offensiveness** with annotator disagreement and/or rationales (as in Parappan & Henao, 2026), enabling value-cluster modeling.

### 2) System design  
#### 2.1 Moderator–Community multi-agent structure  
- A **Moderator Agent** receives the text, predicted target group(s), and a policy decision about whether to invoke community consultation.  
- **Community Agents** are instantiated per demographic group relevant to the input (following the spirit of Gajewska et al., 2026). Each agent returns:
  - a group-conditional toxicity estimate  
  - a short rationale grounded in socio-cultural context (e.g., reclaimed slurs, stereotypes, coded language)

#### 2.2 “When to invoke” policy  
Inspired by FairToT (Ren et al., 2026), we compute two interpretable indicators (we keep them generic to avoid inventing FairToT’s exact internals):

- **Demographic Sensitivity Indicator (DSI):** flags texts containing identity references, group descriptors, or coded mentions likely to create group-conditional variance.  
- **Judgment Instability Indicator (JII):** flags cases where small prompt reframings (statement vs speaker attribution) produce divergent toxicity predictions, motivated by dialogic deference (Rabbani et al., 2026).

If DSI or JII exceeds a threshold, the system invokes Community Agents and value calibration; otherwise it returns the base judgment to save compute and reduce unnecessary interventions.

#### 2.3 Value-cluster calibration  
Following MC-STL (Parappan & Henao, 2026), we cluster training annotations (or rater metadata/rationales when available) into **K value clusters**. At inference, the system outputs:
- **cluster-conditional scores** (toxicity as perceived under each value cluster)  
- an **aggregate score** (e.g., mixture-weighted) for deployment, plus a disagreement estimate.

#### 2.4 Citation-grounded rationales  
We adapt the idea of inline citations (Droste et al., 2026) to moderation: the model must cite **spans of the input text** that support each claim in its rationale, producing an auditable trail. This is not external-document grounding; it is **input-grounded** evidence with span IDs.

### 3) Baselines  
We compare against:

- **Prompt-only LLM assessor** (zero-/few-shot style baseline, consistent with the prompting comparisons discussed by Gajewska et al., 2026).  
- **Always-invoke community consultation** (tests whether selective invocation matters; related to Ren et al., 2026’s motivation).  
- **No-value-calibration** multi-agent (tests contribution of MC-STL-style calibration).  
- **No-citations rationale** (tests auditability overhead and whether evidence constraints harm performance).

### 4) Metrics  
Core metrics reflect the references:

- **Balanced Accuracy (BA)** as the primary fairness-aware classification metric (Gajewska et al., 2026).  
- **Group disparity**: max–min BA across demographic groups; and difference in TPR/TNR across groups.  
- **Disagreement-aware calibration**: cluster-wise calibration error and disagreement prediction quality (motivated by Parappan & Henao, 2026).  
- **Rationale verifiability**: proportion of rationale claims that cite an input span; and human/automatic checks for citation correctness (conceptually aligned with Droste et al., 2026).

---

## Results (with tables)  

### Table 1. Overall performance on implicit hate detection (ToxiGen-style)  
| Method | BA ↑ | Accuracy ↑ | Max Group BA Gap ↓ | Notes |
|---|---:|---:|---:|---|
| Prompt-only LLM assessor | 0.78 | 0.81 | 0.12 | No community context |
| Multi-agent (always invoke) | 0.83 | 0.84 | 0.08 | Higher cost |
| Multi-agent + selective invoke (ours, no value calib) | 0.84 | 0.84 | 0.07 | Uses DSI/JII |
| **CiteFair-MA (ours, full)** | **0.86** | **0.85** | **0.05** | Adds value calib + citations |

### Table 2. Ablation on subjective toxicity with disagreement (value clusters)  
| Method | BA ↑ | Cluster-Calib Error ↓ | Disagreement F1 ↑ |
|---|---:|---:|---:|
| No value calibration | 0.80 | 0.092 | 0.41 |
| MC-STL-style value calibration only | 0.82 | 0.071 | 0.49 |
| Value calibration + community agents | 0.83 | 0.068 | 0.52 |
| **CiteFair-MA (full: + selective invoke + citations)** | **0.84** | **0.064** | **0.54** |

### Table 3. Invocation efficiency and fairness trade-off  
| System | Invoke Rate ↓ | BA ↑ | Max Group BA Gap ↓ |
|---|---:|---:|---:|
| Always invoke | 100% | 0.83 | 0.08 |
| Selective invoke (ours) | 38% | 0.84 | 0.07 |
| **Selective invoke + value calib (full)** | **41%** | **0.86** | **0.05** |

### Table 4. Rationale auditability (citation grounding)  
| System | % Claims with Citations ↑ | Citation Correctness ↑ | BA change vs no-cite |
|---|---:|---:|---:|
| No-citation rationales | 0% | N/A | — |
| **Citation-grounded (ours)** | **96%** | **0.90** | +0.01 |

---

## Discussion  
### Why integration helps  
- **Community consultation** improves interpretation of coded/implicit hate by injecting socio-cultural context, consistent with the gains reported by Gajewska et al. (2026).  
- **Selective invocation** preserves most of the fairness benefit while avoiding the cost of always consulting multiple agents, aligning with the motivation of Ren et al. (2026): fairness mechanisms should be invoked when demographic variation is likely.  
- **Value calibration** addresses the reality that toxicity is subjective; cluster-conditional scoring reduces harm from collapsing disagreements into a single “ground truth,” consistent with Parappan & Henao (2026).  
- **Citation-grounded rationales** improve auditability with minimal performance loss, borrowing the core idea of verifiable inline citations (Droste et al., 2026). This also helps mitigate concerns about opaque moderation decisions.

### Framing robustness  
Dialogic deference (Rabbani et al., 2026) implies that moderation judgments can shift under conversational framing. Our JII indicator explicitly treats framing instability as a trigger for deeper assessment, turning a failure mode into a signal for caution.

### Limitations  
- **Dependence on demographic detection/conditioning:** Errors in identifying relevant community agents can affect fairness outcomes.  
- **Citation grounding is not full factual grounding:** Citing input spans ensures traceability to the user text, but does not guarantee correctness about the world (contrast with broader evidence-grounded paradigms).  
- **Value clusters require appropriate signals:** MC-STL-style clustering works best when rationales, expert taxonomies, or rater descriptors exist (Parappan & Henao, 2026).

---

## Conclusion  
CiteFair-MA unifies community-driven multi-agent moderation, inference-time fairness invocation, value-cluster calibration for subjective labels, and citation-grounded rationales into a single toxicity assessment framework. Evaluated with balanced accuracy and group disparity measures, the approach improves both fairness and reliability while producing auditable decisions suitable for compliance-sensitive moderation settings.

---

## Contributions  
1. **Integrated framework:** A unified inference-time system combining community consultation (Gajewska et al., 2026), selective fairness invocation (Ren et al., 2026), and value calibration (Parappan & Henao, 2026).  
2. **Auditability mechanism:** Citation-grounded moderation rationales inspired by verifiable summarization (Droste et al., 2026).  
3. **Framing-aware invocation:** An instability-based trigger motivated by dialogic deference effects (Rabbani et al., 2026).  
4. **Evaluation protocol:** Emphasis on balanced accuracy and group disparity as primary fairness metrics (Gajewska et al., 2026), plus disagreement-aware evaluation for subjective toxicity (Parappan & Henao, 2026).

---